This section reviews the background that motivates our proposal and the prior
work it builds on. We begin with the landscape of theorem provers, identifying
which of them stand to benefit from our contributions, and then describe the
data structures that will be extended.

\subsection{Theorem Provers}

\paragraph{Automated Theorem Provers.}
\textit{Automated theorem provers} discharge goals without human guidance, and how
much each stands to gain from our work depends on how centrally it relies on
e-graphs. \textit{SMT solvers} such as Z3~\cite{demoura2008z3} and
cvc5~\cite{barbosa2022cvc5} use congruence closure and e-matching as a core
theory solver, so extensions to the e-graph reach them directly.
\textit{Verification languages} such as Dafny~\cite{leino2010dafny},
Viper~\cite{mueller2016viper}, and F*~\cite{swamy2016fstar} compile programs and
their specifications to SMT queries, and so benefit indirectly, through their
solver backend. \textit{Equality-saturation} tools, including
egg~\cite{willsey2021egg} and the inductive provers TheSy~\cite{singher2021thesy}
and CCLemma~\cite{kurashige2024cclemma}, make the e-graph the central
representation for exploring rewrites, and stand to gain the most. By contrast,
\textit{saturation-based} superposition provers such as
Vampire~\cite{kovacs2013vampire} and E~\cite{schulz2013e} organize their search
around term indexing and ordered rewriting rather than a shared e-graph, and are
the least likely to benefit from this line of work.

\paragraph{Interactive Theorem Provers.}
\textit{Interactive theorem provers} let a user steer the proof while the tool
checks each step. Their kernels rest on rich type theories or higher-order
logics, as in Isabelle~\cite{nipkow2002isabelle}, Coq~\cite{bertot2004coq},
Agda~\cite{norell2009agda}, and Lean~\cite{demoura2021lean4}, and this core
proof checking does not use e-graphs, so our extensions have little to offer it.
Their proof automation is another matter: e-graph-based tactics such as Lean's
\texttt{grind}~\cite{leangrind} and Coq's congruence closure~\cite{corbineau2007congruence}
discharge goals by saturating an e-graph with the available facts, and it is
these components, rather than the kernels, that our work stands to improve.

\subsection{Equational Reasoning}

\paragraph{Union-Find.}
A \textit{union-find}, or disjoint-set structure~\cite{tarjan1975union},
maintains an \textit{equivalence relation} over a set of elements using two
operations: \textit{find} returns a canonical representative of the class
containing an element, so that two elements are equivalent exactly when they
share a representative, and \textit{union} merges two classes by making one
representative point to the other. With union by rank and path compression, both
run in near-constant amortized time~\cite{tarjan1984worstcase}. The structure
thus maintains precisely the equivalence relation generated by the equalities
asserted so far, that is, their reflexive, symmetric, and transitive closure, but
records nothing about the structure of the terms that name the elements.

\paragraph{E-Graphs.}
An \textit{e-graph}~\cite{nelson1980congruence, willsey2021egg} lifts this idea
from opaque elements to functional \textit{terms}, maintaining not merely an equivalence
relation but a \textit{congruent} one: a relation additionally closed under
congruence, so that whenever the arguments of a function are pairwise
equivalent, the applications are too. A term is stored as an \textit{e-node}, a
function symbol applied to a tuple of child \textit{e-classes}, and an e-class is
a set of e-nodes that have been proven equal. Internally, an e-graph is a
union-find over e-class identifiers together with a hash-cons mapping each e-node
to its e-class, so \textit{find} and \textit{union} carry over unchanged.
Merging two e-classes with \textit{union}, however, can break the congruence
invariant, since two e-nodes that have just become congruent may still sit in
different e-classes. For example, suppose an e-graph holds the terms $f(a)$ and
$f(b)$ in distinct e-classes; asserting $a = b$ merges the classes of $a$ and
$b$, so congruence now demands $f(a) = f(b)$, yet the two applications still sit
in separate e-classes until the invariant is restored. The \textit{rebuild}
operation restores the invariant by
repeatedly merging such newly congruent e-nodes until a fixpoint is
reached~\cite{willsey2021egg}. An e-graph therefore represents the
\textit{congruent equivalence relation} generated by the asserted equalities,
that is, their reflexive, symmetric, transitive, and congruent
closure~\cite{zakhour2025disequality}.

\paragraph{Equality Saturation.} \textit{Equality saturation}~\cite{tate2009eqsat, willsey2021egg}
is a framework built around e-graphs that repeatedly applies a set of rewrite 
rules, non-destructively so that
rules need not be ordered, until it is \textit{saturated} or a resource bound is
reached. Then, one can \textit{extract} an optimal representative under a cost
model or prove two terms equal under the axioms of the rewrite rules. In detail,
each rewrite $\ell \to r$ is applied by \textit{e-matching}, finding all
instances of the pattern $\ell$ up to the equivalences the e-graph
records~\cite{demoura2007ematching} and making $r$ equal to every match.
Sidestepping the \textit{phase-ordering} problem this way, equality saturation
underlies applications from compiler optimization~\cite{tate2009eqsat} to
inductive theorem proving, as in TheSy~\cite{singher2021thesy} and
CCLemma~\cite{kurashige2024cclemma}, while the same e-matching mechanism drives
quantifier instantiation in SMT solvers~\cite{demoura2007ematching}.